% concatenated from main.tex
\documentclass[12pt,a4paper]{amsart}
\usepackage{amsmath,amssymb,amsthm}
\usepackage{geometry}
\usepackage{hyperref}

\allowdisplaybreaks[4]
\geometry{margin=2.5cm}

%------------------------------------------------
% Theorem environments
%------------------------------------------------

\newtheorem{theorem}{Theorem}[section]
\newtheorem{lemma}[theorem]{Lemma}
\newtheorem{proposition}[theorem]{Proposition}
\newtheorem{corollary}[theorem]{Corollary}

\theoremstyle{definition}
\newtheorem{definition}[theorem]{Definition}
\newtheorem{remark}[theorem]{Remark}

%------------------------------------------------
% Macros
%------------------------------------------------

\newcommand{\Z}{\mathbb Z}

%------------------------------------------------
% Title
%------------------------------------------------

\title[]{Admissible Bases and Generic Lattice Ideals}
\author[]{Anargyros Katsabekis}
\address {Department of Mathematics, University of Ioannina, 45110 Ioannina, Greece} \email{katsampekis@uoi.gr}
\keywords{Positive lattice; Generic lattice ideal; Admissible basis;
Neighbor of the origin; Gr\"obner basis; Algebraic Scarf complex}
\subjclass{13F65, 13F20, 13P10, 13D02}

\begin{document}

\begin{abstract} We introduce the notion of an admissible basis for rank-three positive
lattices in $\mathbb Z^4$, namely a $\mathbb Z$-basis
$\{\mathbf c_1,\mathbf c_2,\mathbf c_3\}$ satisfying three explicit
sign conditions~\textup{(I)}, \textup{(II)}, and~\textup{(III)} on the
coordinates of the basis vectors. We study lattice vectors of the form
$
\mathbf u^{(\mu,\lambda)}
=
\mu\mathbf c_1+\mathbf c_2+\lambda\mathbf c_3,$
where $\mu$ and $\lambda$ are positive integers. Under
conditions~\textup{(I)}, \textup{(II)}, and~\textup{(III)}, we obtain a
complete characterization of the vectors
$\mathbf u^{(\mu,\lambda)}$ with positive first and fourth coordinates
that are neighbors of the origin: such a vector is a neighbor if and
only if $\mu=1$ and $\lambda\le2$.

As an application, for every integer $m\ge1$ we construct a positive
lattice $L(m)\subseteq\mathbb Z^4$ admitting an admissible basis whose
associated lattice ideal is generic. This yields an infinite family of
generic lattice ideals with exactly seven minimal binomial generators. This family shows that the characterization is sharp, since both
$\mathbf u^{(1,1)}$ and $\mathbf u^{(1,2)}$ occur as neighbors of the
origin.
\end{abstract}

\maketitle

\section{Introduction}

Lattice ideals form a fundamental class of binomial ideals arising in
combinatorics, commutative algebra, and algebraic geometry. They
include, in particular, the defining ideals of affine monomial curves.
If an affine monomial curve is parametrized by
$x_1=t^{n_1},\ldots,x_r=t^{n_r}$, where
$\gcd(n_1,\ldots,n_r)=1$, then its defining ideal is the lattice ideal
associated with the lattice
\[
\ker_{\mathbb Z}(n_1,\ldots,n_r)
=\{(u_1,\ldots,u_r)\in\mathbb Z^r \mid
u_1n_1+\cdots+u_rn_r=0\}.
\]
More generally, we consider lattice ideals
associated with positive lattices, namely lattices
$L\subseteq\mathbb Z^r$ satisfying
$
L\cap\mathbb N^r=\{\mathbf0\}.$ Throughout the paper, $K$ denotes a field.

An important problem in the theory of lattice ideals is the explicit
determination of the minimal free resolution of the quotient ring
$K[x_1,\ldots,x_r]/I_L$,
where $I_L$ denotes the lattice ideal associated with a positive
lattice $L$. For generic lattice ideals this problem is solved by
\cite[Theorem~4.2]{PS}, which shows that the minimal free resolution is
supported on the algebraic Scarf complex. Moreover, the minimal binomial generators of $I_L$ are in bijection with the neighbors of the origin; see Proposition~\ref{cor:neighbors}. Thus, determining the neighbor set provides direct information about the minimal generators
of $I_L$. This motivates the problem of determining the neighbor set
directly from the lattice, without first computing the associated
lattice ideal or a Gr\"obner basis.

Determining the neighbor set directly from a lattice is, however, a
difficult problem. Existing approaches typically proceed by first
computing the associated lattice ideal, for example via a Gr\"obner
basis or a minimal binomial generating set. As a first step, we prove
a coordinatewise characterization of the neighbors of the origin for
arbitrary positive lattices; see Lemma~\ref{lem:neighbor}. Building on
this criterion, we introduce admissible bases for positive lattices of
rank three in $\mathbb Z^4$ and use them to obtain a complete
characterization of the vectors
$
\mu\mathbf c_1+\mathbf c_2+\lambda\mathbf c_3$
having positive first and fourth coordinates that are neighbors of the
origin. We also construct an explicit infinite family of generic
lattice ideals for which the reduced Gr\"obner bases, minimal binomial
generating sets, neighbor sets, Betti numbers, and the $f$-vectors of
the associated algebraic Scarf complexes are determined explicitly.

In the study of hull resolutions of generic affine monomial curves,
Ojeda and Pis\'on Casares~\cite{OP}, building on work of
Barany and Scarf~\cite{BS}, considered lattice vectors of the form
$\mathbf c_1+\mathbf c_2+\lambda\mathbf c_3$, where
$\mathbf c_1,\mathbf c_2,\mathbf c_3$ are basis elements of $L$
satisfying certain sign conditions. Such vectors arise naturally as
neighbors of the origin; see Section~4 of~\cite{OP}. This raises the question of which vectors of the form
$\mathbf c_1+\mathbf c_2+\lambda\mathbf c_3$
can occur as neighbors of the origin and, more generally, which vectors
of the form
$\mathbf u^{(\mu,\lambda)}
=\mu\mathbf c_1+\mathbf c_2+\lambda\mathbf c_3$,
where $\mu$ and $\lambda$ are positive integers, can be neighbors.
The present paper addresses this question by establishing explicit
criteria, in terms of an admissible basis, that restrict when vectors
of this form can occur as neighbors of the origin.

To make this precise, we introduce the notion of an \emph{admissible
basis} $\{\mathbf c_1,\mathbf c_2,\mathbf c_3\}$ of a rank-three
lattice $L\subseteq\mathbb Z^4$: a $\mathbb Z$-basis satisfying three
explicit sign conditions~\textup{(I)}, \textup{(II)}, and
\textup{(III)} on the entries of the basis vectors
(see Definition~\ref{def:setup}).
The three conditions are motivated by Lemma~\ref{lem:neighbor} and are
precisely those needed to apply its coordinatewise criterion in
Section~\ref{sec:lambda}. Conditions~\textup{(I)} and~\textup{(II)}
isolate precisely the coordinate inequalities needed in the proofs of
Section~\ref{sec:lambda}: condition~\textup{(I)} is used to compare
first coordinates, while condition~\textup{(II)} controls the second
coordinates. Condition~\textup{(III)} is satisfied by the explicit
family constructed in Section~\ref{sec:setup} and guarantees the
required positivity properties of the fourth coordinates. Thus an
admissible basis records exactly the lattice-theoretic information
needed for the arguments based on Lemma~\ref{lem:neighbor}, without
computing the associated lattice ideal or a Gröbner basis.

The main structural results concern vectors
$\mathbf u^{(\mu,\lambda)}$
having positive first and fourth coordinates,
where this notion is defined in
Section~\ref{sec:lambda}.
For every positive lattice admitting an admissible basis, we obtain a
complete characterization of the vectors of this form that are
neighbors of the origin: such a vector is a neighbor if and only if
$\mu=1$ and $\lambda\le2$.
Moreover, both values
$\lambda=1$ and $\lambda=2$
occur in the explicit family constructed in
Section~\ref{sec:setup}, showing that the characterization is sharp.

For generic lattice ideals, the neighbors of the origin are in
bijection with the minimal binomial generators.
Thus the above characterization provides an explicit
lattice-theoretic criterion for determining a large class of minimal
generators directly from an admissible basis, without computing the
associated lattice ideal or a Gr\"obner basis.

The second main contribution of the paper is the construction and
analysis of an explicit infinite family of generic lattice ideals
having exactly seven minimal binomial generators.
Although isolated examples of generic lattice ideals with seven minimal
binomial generators were already known, no infinite family of such
ideals appears to have been constructed previously.
In contrast, Ojeda~\cite{Ojeda} proved that for every integer
$m\ge8$ there exists a generic lattice ideal of codimension three
having exactly $m$ minimal binomial generators.
Our construction complements the existence theorem of
Ojeda~\cite{Ojeda} by providing an explicit infinite family of generic
lattice ideals having exactly seven minimal binomial generators.

More precisely, for every integer $m\ge1$ we construct a positive
lattice $L(m)$ admitting an admissible basis such that the associated
lattice ideal $I_{L(m)}$ is generic.
For every member of the family we determine explicitly the reduced
Gr\"obner basis, the unique minimal binomial generating set, the
neighbor set, the Betti numbers of $K[x_1,\ldots,x_4]/I_{L(m)}$,
and the $f$-vector of the associated algebraic Scarf complex.
If $5$ does not divide $m$, then $\mathbb Z^4/L(m)$ is torsion free, so
$I_{L(m)}$ is the defining ideal of an affine monomial curve.
Furthermore, the family shows that the bound $\lambda\le2$ in
Theorem~\ref{thm:mu-lambda} is sharp, since both values
$\lambda=1$ and $\lambda=2$ occur.

The paper is organized as follows.
Section~\ref{sec:preliminaries} introduces the notation and basic
results used throughout the paper. Its principal new result is
Lemma~\ref{lem:neighbor}, which provides a coordinatewise
characterization of the neighbors of the origin for arbitrary positive
lattices and serves as the main tool in the subsequent sections.

Section~\ref{sec:setup} introduces admissible bases and constructs an
explicit infinite family of positive lattices admitting such bases.
Theorem~\ref{thm:groebner} shows that the associated lattice ideals are
generic with exactly seven minimal binomial generators, determines
their reduced Gr\"obner bases and neighbor sets, and proves that the
Gr\"obner basis is the unique minimal binomial generating set.
Proposition~\ref{prop:scarf} determines the Betti numbers and the
$f$-vector of the associated algebraic Scarf complex.

Section~\ref{sec:lambda} develops a lattice-theoretic approach to the
study of neighbors of the origin based on admissible bases. It
culminates in Corollary~\ref{cor:classification}, which gives a
complete characterization of the vectors
$\mathbf u^{(\mu,\lambda)}$ having positive first and fourth
coordinates that are neighbors of the origin: such a vector is a
neighbor of the origin if and only if $\mu=1$ and $\lambda\le2$.




%--------------------------------------------------------------------
\section{Preliminaries}
\label{sec:preliminaries}
%--------------------------------------------------------------------

In this section we recall the notions and results concerning lattice
ideals, Gr\"obner bases, and algebraic Scarf complexes that will be
used throughout the paper. Standard references are
\cite{CLO,MS,PS,Stu}.

Let $L\subseteq\mathbb Z^{r}$ be a nonzero positive lattice, that is,
\[
L\cap\mathbb N^{r}=\{\mathbf 0\}.
\] For a vector
${\bf u}=(u_1,\ldots,u_r)\in\mathbb Z^{r}$,
we denote its $i$-th coordinate by
$({\bf u})_i=u_i$, for every $1\le i\leq r$.
Its \emph{support} is the set
\[
\operatorname{supp}({\bf u})
=
\{\,i\in\{1,\ldots,r\}\mid({\bf u})_i\neq0\,\}.
\]
We say that ${\bf u}$ has \emph{full support} if
\[
\operatorname{supp}({\bf u})=\{1,\ldots,r\}.
\]

Every vector ${\bf u}\in\mathbb Z^{r}$ admits a unique decomposition
\[
{\bf u}={\bf u}^{+}-{\bf u}^{-},
\]
where the $i$-th coordinates of ${\bf u}^{+}$ and ${\bf u}^{-}$ are
given, for $i=1,\ldots,r$, by
\[
({\bf u}^{+})_i=\max\{u_i,0\},\ \textrm{and} \
({\bf u}^{-})_i=\max\{-u_i,0\}.
\]
The vectors ${\bf u}^{+}$ and ${\bf u}^{-}$ are called the
\emph{positive part} and the \emph{negative part} of ${\bf u}$,
respectively. Their coordinates are nonnegative, and they have
disjoint supports.

For vectors ${\bf a},{\bf b}\in\mathbb Z^{r}$, we write
$
{\bf a}\le{\bf b}$
if
$
({\bf a})_i\le({\bf b})_i$ for every $1 \leq i \leq r.$

Let $S=K[x_1,\ldots,x_r]$
be the polynomial ring in $r$ variables over $K$.
For
${\bf a}=(a_1,\ldots,a_r)\in\mathbb N^{r}$,
we write
$
{\bf x}^{\bf a}=x_1^{a_1} \cdots x_r^{a_r}.$ The lattice $L$ determines the ideal
\[
I_L=
\left\langle
{\bf x}^{{\bf u}^{+}}-{\bf x}^{{\bf u}^{-}}
\mid
{\bf u}\in L
\right\rangle
\subseteq S,
\]
called the \emph{lattice ideal} associated with $L$.
Since $L$ is positive, $I_L$ is homogeneous with respect to a positive
grading~\cite{PS}. Each generator
$
{\bf x}^{{\bf u}^{+}}-{\bf x}^{{\bf u}^{-}}
$
is the difference of two monomials, and hence is a binomial.
Therefore $I_L$ is a binomial ideal.
For ${\bf u}\in L$, we call
$
{\bf x}^{{\bf u}^{+}}-{\bf x}^{{\bf u}^{-}}
$
the \emph{binomial associated with} ${\bf u}$.
Its support is the support of ${\bf u}$.
If the quotient group
$\mathbb Z^{r}/L$
is torsion free, then $I_L$ is prime; in this case it is the toric
ideal associated with $L$.

A lattice ideal $I_L$ is called \emph{generic} if it is generated by
binomials of full support. Generic lattice ideals play a distinguished role in the theory of
lattice ideals because, by results of Peeva and
Sturmfels~\cite{PS}, their minimal binomial generators are determined
by the neighbors of the origin, and their minimal free resolutions are
supported on the algebraic Scarf complex.

A central role in the study of generic lattice ideals is played by the
neighbors of the origin. For a finite subset
$J=\{\mathbf u^{(1)},\ldots,\mathbf u^{(m)}\}\subseteq L$,
where
$\mathbf u^{(k)}=(u^{(k)}_1,\ldots,u^{(k)}_r)$ for $1 \leq k \leq m$,
define the componentwise maximum of $J$ by
\[
\max(J)
=
\bigl(
\max_{1\le k\le m}u^{(k)}_1,\,
\max_{1\le k\le m}u^{(k)}_2,\,
\ldots,
\max_{1\le k\le m}u^{(k)}_r
\bigr).
\] Following~\cite[Section~2]{PS},
let $\Delta_L$ be the simplicial complex on the lattice $L$ defined by
\[
\Delta_L
=
\left\{
J\subseteq L:
\max(J)\neq\max(J')
\text{ for every subset }
J'\subseteq L
\text{ with }
J'\neq J
\right\}.
\]
Let
\[
\Delta_L^0
=
\{J\subseteq L\setminus\{\mathbf0\}:J\cup\{\mathbf0\}\in\Delta_L\}.
\]
The vertices of $\Delta_L^0$ are called the
\emph{neighbors of the origin}. We denote the set of all neighbors of
the origin by $N(L)$.

The following lemma provides a coordinatewise criterion for determining whether a lattice vector is a neighbor of the origin. The necessity is implicit in the proof of Proposition~2.2 of \cite{PS}; we include a short proof for completeness and establish the converse.

\begin{lemma}\label{lem:neighbor}
Let $L\subseteq\mathbb Z^r$ be a positive lattice and let
$\mathbf u\in L\setminus\{\mathbf0\}$.
Then $\mathbf u$ is a neighbor of the origin if and only if there is
no $\mathbf b\in L\setminus\{\mathbf0,\mathbf u\}$ such that
$\mathbf b\le\mathbf u^+$.
\end{lemma}
\begin{proof}
Suppose that $\mathbf u$ is a neighbor of the origin. Then
$\{\mathbf u\}\in\Delta_L^0$, and therefore
$\{\mathbf0,\mathbf u\}\in\Delta_L$. Assume that there exists
$\mathbf b\in L\setminus\{\mathbf0,\mathbf u\}$
such that
$\mathbf b\le\mathbf u^+$.
For each $1 \leq i \leq r$, we have
$
(\mathbf0)_i=0\le(\mathbf u^+)_i,$ 
$(\mathbf u)_i\le(\mathbf u^+)_i,$ and
$(\mathbf b)_i\le(\mathbf u^+)_i.$
Since
$(\mathbf u^+)_i=\max\{0,(\mathbf u)_i\}$, either $(\mathbf0)_i$ or $(\mathbf u)_i$ is equal to
$(\mathbf u^+)_i$. Hence $
\max\{(\mathbf0)_i,(\mathbf u)_i,(\mathbf b)_i\}
=
(\mathbf u^+)_i.$
Thus
\[
\max(\{\mathbf0,\mathbf u,\mathbf b\})
=
\mathbf u^+
=
\max(\{\mathbf0,\mathbf u\}).
\]
Since
$\{\mathbf0,\mathbf u,\mathbf b\}$ and
$\{\mathbf0,\mathbf u\}$
are distinct subsets of $L$ with the same componentwise maximum, this
contradicts the fact that
$\{\mathbf0,\mathbf u\}\in\Delta_L$.
Therefore no such vector $\mathbf b$ exists.

Conversely, suppose that there is no
$\mathbf b\in L\setminus\{\mathbf0,\mathbf u\}$
satisfying
$\mathbf b\le\mathbf u^+$.
Since $L$ is positive and $\mathbf u\neq\mathbf0$, neither
$\mathbf u$ nor $-\mathbf u$ belongs to $\mathbb N^r$.
Hence
$\mathbf u\neq\mathbf u^+$.
Suppose that
$\{\mathbf0,\mathbf u\}\notin\Delta_L$.
Since
$
\max(\{\mathbf0,\mathbf u\})
=
\mathbf u^+,$
there exists a finite subset
$J\subseteq L$
such that
$J\neq\{\mathbf0,\mathbf u\}$
and
$\max(J)=\mathbf u^+$. Thus every element of $J$ is bounded coordinatewise by
$\mathbf u^+$; that is,
$\mathbf c\le\mathbf u^+$ for every $\mathbf c\in J$.
Since
$J\neq\{\mathbf0,\mathbf u\}$,
there exists
$\mathbf b\in L\setminus\{\mathbf0,\mathbf u\}$
such that
$\mathbf b\in J$.
Therefore,
$\mathbf b\le\mathbf u^+$,
contradicting the hypothesis.
Hence
$\{\mathbf0,\mathbf u\}\in\Delta_L$.
Equivalently,
$\{\mathbf u\}\in\Delta_L^0$,
so $\mathbf u$ is a neighbor of the origin.
\end{proof}
The following proposition gives the correspondence between the
neighbors of the origin and the minimal binomial generators of a
generic lattice ideal.

\begin{proposition} \label{cor:neighbors}
Let $I_L$ be a generic lattice ideal. A lattice vector
${\bf u}\in L$ is a neighbor of the origin if and only if
$
{\bf x}^{{\bf u}^{+}}-{\bf x}^{{\bf u}^{-}}$
is a minimal binomial generator of $I_L$.
\end{proposition}

\begin{proof}
By~\cite[Theorem~2.2]{Sabzrou}, for a generic lattice ideal $I_L$,
${\bf u}\in L$ is a neighbor of the origin if and only if
$
B={\bf x}^{{\bf u}^{+}}-{\bf x}^{{\bf u}^{-}}$
is an indispensable binomial of $I_L$, that is, either $B$ or $-B$
belongs to every binomial generating set of $I_L$.
By~\cite[Remark~4.4(3)]{PS}, a generic lattice ideal has a unique
minimal binomial generating set. Hence the minimal binomial generators
coincide with the indispensable binomials. Therefore
$
{\bf x}^{{\bf u}^{+}}-{\bf x}^{{\bf u}^{-}}$ is a minimal binomial generator of $I_L$ if and only if
${\bf u}$ is a neighbor of the origin.
\end{proof}
We shall also use standard facts from the theory of Gr\"obner bases.
Fix a monomial order $\prec$ on $S$. For a polynomial $f\in S$, we
denote by $\operatorname{lm}(f)$ and $\operatorname{lc}(f)$ its
leading monomial and leading coefficient, respectively, with respect
to $\prec$. Our computations rely on Buchberger's criterion together
with the coprimeness criterion.

\begin{theorem}[Buchberger's Criterion {\cite[Chapter~2, \S6, Theorem~6]{CLO}}]
\label{thm:buchberger}
A finite subset $G=\{g_1,\ldots,g_t\}\subseteq I$
is a Gr\"obner basis of an ideal $I\subseteq S$ if and only if every
$S$-polynomial $S(g_i,g_j)$ reduces to zero modulo $G$.
\end{theorem}

\begin{proposition}[{\cite[Chapter~2, Proposition~4]{CLO}}]
\label{prop:coprime}
Let $f,g\in S$. If the leading monomials of $f$ and $g$ are relatively
prime, then the $S$-polynomial $S(f,g)$ reduces to zero modulo
$\{f,g\}$.
\end{proposition}

%--------------------------------------------------------------------
\section{Admissible Lattices and the Parametric Family}
\label{sec:setup}
%--------------------------------------------------------------------
In this section we first introduce the notion of an admissible basis
for a rank-three lattice in $\mathbb Z^4$. The defining sign
conditions ensure that the lattice vectors introduced below have full
support and satisfy the coordinate properties needed in the subsequent
analysis of neighbors of the origin.

\begin{definition}\label{def:setup} A $\mathbb Z$-basis
$\{\mathbf c_1,\mathbf c_2,\mathbf c_3\}$ of a rank-three lattice
$L\subseteq\mathbb Z^4$ is called \emph{admissible} if $\mathbf c_1=(\alpha_{11},-\alpha_{12},-\alpha_{13},-\alpha_{14})$,
$\mathbf c_2=(-\alpha_{21},-\alpha_{22},-\alpha_{23},\alpha_{24})$, and
$\mathbf c_3=(-\alpha_{31},\alpha_{32},-\alpha_{33},-\alpha_{34})$,
where every $\alpha_{ij}$ is a positive integer, and the following conditions hold:
\begin{itemize}
\item[\textbf{(I)}]
$\alpha_{11}>\alpha_{21}+2\alpha_{31}$,
\item[\textbf{(II)}]
$\max(\alpha_{12},\alpha_{22})<\alpha_{32}
<\alpha_{12}+\alpha_{22}$,
\item[\textbf{(III)}]
$\alpha_{24}>\alpha_{14}+2\alpha_{34}$.
\end{itemize}

\end{definition}

The inequalities in Definition~\ref{def:setup} are chosen so that the
coordinatewise criterion of Lemma~\ref{lem:neighbor} can be applied
effectively to the vectors studied in
Section~\ref{sec:lambda}.

Throughout the remainder of the paper, given an admissible basis, we
write $\mathbf u_i=\mathbf c_i$ for $i=1,2,3$, and define
\[
\mathbf u_4=\mathbf c_1+\mathbf c_3,\;
\mathbf u_5=\mathbf c_2+\mathbf c_3,\;
\mathbf u_6=\mathbf c_1+\mathbf c_2+\mathbf c_3,\;
\mathbf u_7=\mathbf c_1+\mathbf c_2+2\mathbf c_3.
\]

The following proposition establishes a basic consequence of the
defining properties of an admissible basis that will be used
repeatedly throughout the paper.

\begin{proposition}\label{prop:full-support} Under conditions
\textup{(I)}, \textup{(II)}, and \textup{(III)}, every vector $\mathbf u_i$, $1\le i\le7$, has full support. Moreover, $(\mathbf u_i^+)_3=0$ and $(\mathbf u_i^-)_3>0$ for every $1 \leq i \leq 7$.
\end{proposition}

\begin{proof}
The third coordinates of the vectors
$\mathbf u_1,\ldots,\mathbf u_7$ are
\[
-\alpha_{13},\,
-\alpha_{23},\,
-\alpha_{33},\,
-(\alpha_{13}+\alpha_{33}),\,
-(\alpha_{23}+\alpha_{33}),\,
-(\alpha_{13}+\alpha_{23}+\alpha_{33}),\,
-(\alpha_{13}+\alpha_{23}+2\alpha_{33}).
\]
Since every $\alpha_{ij}$ is positive, each of these numbers is
negative. Hence
$
(\mathbf u_i^+)_3=0$ and $(\mathbf u_i^-)_3>0$ for every $1\le i\le7$.

It remains to verify that every $\mathbf u_i$ has full support.
The first coordinates of the vectors
$\mathbf u_1,\ldots,\mathbf u_7$ are
\[
\alpha_{11},\,
-\alpha_{21},\,
-\alpha_{31},\,
\alpha_{11}-\alpha_{31},\,
-(\alpha_{21}+\alpha_{31}),\,
\alpha_{11}-\alpha_{21}-\alpha_{31},\,
\alpha_{11}-\alpha_{21}-2\alpha_{31}.
\]
Since
$
\alpha_{11}>\alpha_{21}+2\alpha_{31}>0$
by condition~\textup{(I)}, the last expression is positive, while the
remaining quantities are clearly nonzero. Hence every first coordinate
is nonzero.

The second coordinates of the vectors
$\mathbf u_1,\ldots,\mathbf u_7$ are
\[
-\alpha_{12},\,
-\alpha_{22},\,
\alpha_{32},\,
\alpha_{32}-\alpha_{12},\,
\alpha_{32}-\alpha_{22},\,
\alpha_{32}-\alpha_{12}-\alpha_{22},\,
2\alpha_{32}-\alpha_{12}-\alpha_{22}.
\]
Condition~\textup{(II)} implies that
$\alpha_{32}-\alpha_{12}>0$, $\alpha_{32}-\alpha_{22}>0$, and 
$\alpha_{32}-\alpha_{12}-\alpha_{22}<0.$
Moreover,
\[
2\alpha_{32}-\alpha_{12}-\alpha_{22}
=
(\alpha_{32}-\alpha_{12})
+
(\alpha_{32}-\alpha_{22})
>
0.
\]
Hence each of these quantities is nonzero.

Finally, the fourth coordinates are
\[
-\alpha_{14},\,
\alpha_{24},\,
-\alpha_{34},\,
-(\alpha_{14}+\alpha_{34}),\,
\alpha_{24}-\alpha_{34},\,
\alpha_{24}-\alpha_{14}-\alpha_{34},\,
\alpha_{24}-\alpha_{14}-2\alpha_{34}.
\]
Since $ \alpha_{24}>\alpha_{14}+2\alpha_{34}$ by condition ~\textup{(III)}, the last three quantities are positive, while the first four are clearly nonzero. Hence every fourth coordinate is nonzero. Therefore each $\mathbf u_i$
has full support.
\end{proof}

Next we introduce an explicit parametric family of admissible lattices that will serve as the main source of examples throughout the remainder of the paper.

\begin{definition}\label{def:family}
For each integer $m\ge1$, let $\mathbf c_1(m)=(m+4,-(m+1),-1,-1)$,
$\mathbf c_2(m)=(-(m+1),-3,-1,m+4)$, and
$\mathbf c_3(m)=(-1,m+3,-(m+1),-1)$. Set
\[
L(m)=
\mathbb Z\mathbf c_1(m)+
\mathbb Z\mathbf c_2(m)+
\mathbb Z\mathbf c_3(m).
\]
\end{definition}


The following proposition establishes the basic properties of this
family.

\begin{proposition}\label{prop:family}
For every integer $m\ge1$, let
\[
\mathbf n(m)=
\bigl(n_1(m),n_2(m),n_3(m),n_4(m)\bigr),
\]
where
\[
\begin{aligned}
n_1(m)&=m^3+7m^2+21m+20, &
n_2(m)&=m^3+8m^2+25m+25,\\
n_3(m)&=m^3+8m^2+26m+25, &
n_4(m)&=m^3+8m^2+28m+30.
\end{aligned}
\]
Then the following hold.

\begin{enumerate}
\item[\textup{(i)}]
\[
\gcd(n_1(m),n_2(m),n_3(m),n_4(m))
=
\begin{cases}
5,&\text{if $5$ divides $m$,}\\
1,&\text{otherwise.}
\end{cases}
\]

\item[\textup{(ii)}]
The lattice $L(m)$ has admissible basis
$\{\mathbf c_1(m),\mathbf c_2(m),\mathbf c_3(m)\}$.

\item[\textup{(iii)}]
If $5$ does not divide $m$, then $
L(m)={\rm ker}_{\mathbb Z}(n_{1}(m),n_{2}(m),n_{3}(m),n_{4}(m)).$
\end{enumerate}
\end{proposition}

\begin{proof}
Let $A(m)$ be the $3\times4$ integer matrix whose rows are
$\mathbf c_1(m)$, $\mathbf c_2(m)$, and $\mathbf c_3(m)$.
For $i=1,\ldots,4$, let $\Delta_i$ denote the determinant of the
submatrix obtained by deleting the $i$-th column of $A(m)$.
A direct computation yields
\[
(-\Delta_1,\Delta_2,-\Delta_3,\Delta_4)
=
(n_1(m),n_2(m),n_3(m),n_4(m)).
\]
Hence
$
A(m)\mathbf n(m)^{\mathrm T}=\mathbf 0,$ where $\mathbf n(m)^{\mathrm T}$ denotes the transpose of
$\mathbf n(m)$.

We first prove assertion~\textup{(i)}.
Since $n_3(m)-n_2(m)=m$ and $n_4(m)-n_3(m)=2m+5$,
any common divisor of
$n_1(m),n_2(m),n_3(m),n_4(m)$
divides both $m$ and $2m+5$.
Hence
$
\gcd(n_1(m),n_2(m),n_3(m),n_4(m))$
divides $\gcd(m,2m+5)
=
\gcd(m,5).$

If $5$ does not divide $m$, then $\gcd(m,5)=1$, and therefore
\[
\gcd(n_1(m),n_2(m),n_3(m),n_4(m))=1.
\]

Assume now that $5$ divides $m$. Since every nonconstant term of $n_i(m)$ contains a factor of $m$, and the constant term is divisible by $5$, it follows that $5$ divides $n_{i}(m)$ for each $1 \leq i \leq 4$.
Hence $5$ divides
$\gcd(n_1(m),n_2(m),n_3(m),n_4(m))$.
On the other hand,
$
\gcd(n_1(m),n_2(m),n_3(m),n_4(m))$ divides
$\gcd(m,5)=5,$ so
\[
\gcd(n_1(m),n_2(m),n_3(m),n_4(m))=5.
\]

Next we prove assertion~\textup{(ii)}.
The relevant coefficients are
\[
\begin{aligned}
\alpha_{11}&=m+4,\;
\alpha_{12}=\alpha_{21}=m+1,\;
\alpha_{22}=3,\;
\alpha_{31}=1,\\
\alpha_{32}&=m+3,\;
\alpha_{14}=\alpha_{34}=1,\;
\alpha_{24}=m+4.
\end{aligned}
\]

Condition~\textup{(I)} follows from
\[
\alpha_{11}-\alpha_{21}-2\alpha_{31}
=(m+4)-(m+1)-2=1>0.
\]

For condition~\textup{(II)},
\[
\alpha_{12}+\alpha_{22}
=(m+1)+3=m+4,
\]
and therefore
\[
\alpha_{32}=m+3<m+4=\alpha_{12}+\alpha_{22}.
\]
Moreover,
\[
\max\{\alpha_{12},\alpha_{22}\}
=\max\{m+1,3\}
<m+3=\alpha_{32}.
\]

Finally,
\[
\alpha_{24}-\alpha_{14}-2\alpha_{34}
=(m+4)-1-2=m+1>0,
\]
so condition~\textup{(III)} is satisfied.

Finally, we prove assertion~\textup{(iii)}.
Since $A(m)$ has rank~$3$, its Smith normal form is
${\rm diag}(e_1,e_2,e_3,0),$ where $e_1,e_2,e_3$ are the invariant factors.
The discussion following Theorem~II.9 in
\cite{Newman} shows that
$
e_1e_2e_3
=
\gcd(n_1(m),n_2(m),n_3(m),n_4(m)).$

If $5$ does not divide $m$, then assertion~\textup{(i)} gives
$e_1e_2e_3=1$.
Since $e_1$ divides $e_2$ and $e_2$ divides $e_3$, we obtain
$e_1=e_2=e_3=1$.
By the fundamental theorem of finitely generated abelian groups
(\cite[Section~12.1]{DF}),
\[
\mathbb Z^4/L(m)
\cong
\mathbb Z
\oplus
\mathbb Z/e_1\mathbb Z
\oplus
\mathbb Z/e_2\mathbb Z
\oplus
\mathbb Z/e_3\mathbb Z,
\]
and therefore $\mathbb Z^4/L(m)$ is torsion free.
Moreover,
$
A(m)\mathbf n(m)^{\mathrm T}={\bf 0},$
so
$$
L(m) \subseteq {\rm ker}_{\mathbb Z}(n_{1}(m),n_{2}(m),n_{3}(m),n_{4}(m)).$$
Suppose that $L(m) \neq {\rm ker}_{\mathbb Z}(n_{1}(m),n_{2}(m),n_{3}(m),n_{4}(m))$. Since both lattices have
rank~$3$, the quotient
${\rm ker}_{\mathbb Z}(n_{1}(m),n_{2}(m),n_{3}(m),n_{4}(m))/L(m)$
is a nontrivial finite subgroup of $\mathbb Z^4/L(m)$. Since every element of a finite group has finite order, it follows that $\mathbb Z^4/L(m)$ contains a nontrivial torsion element, contradicting the fact that
$\mathbb Z^4/L(m)$ is torsion-free. Hence
$
L(m)={\rm ker}_{\mathbb Z}(n_{1}(m),n_{2}(m),n_{3}(m),n_{4}(m)).$
\end{proof}

The next proposition establishes that the lattices $L(m)$ are
positive.

\begin{proposition}
\label{prop:positive}
For every integer $m\ge1$, the lattice $L(m)$ is positive.
\end{proposition}

\begin{proof}
Let $\mathbf u=(u_1,\ldots,u_4)\in L(m)\cap\mathbb N^4$. Since
$A(m)\mathbf n(m)^{\mathrm T}={\bf 0}$, the vector $\mathbf n(m)$ is
orthogonal to every vector of $L(m)$. Hence
$
\mathbf n(m)\cdot\mathbf u=0.$
Since each coordinate of $\mathbf n(m)$ is positive and each
$u_i\ge0$, we have
\[
\mathbf n(m)\cdot\mathbf u
=
n_1(m)u_1+\cdots+n_4(m)u_4,
\]
where every summand is nonnegative. Therefore
$\mathbf n(m)\cdot\mathbf u=0$ if and only if
$u_1=u_2=u_3=u_4=0$. Hence
$\mathbf u=\mathbf0$, and consequently
$
L(m)\cap\mathbb N^4=\{\mathbf0\}.$
Thus $L(m)$ is positive.
\end{proof}

\begin{remark}
For $m=1$, Proposition~\ref{prop:family} gives
$\mathbf n(1)=(49,59,60,67).$ The affine monomial curve parametrized by
$x_1=t^{49}$, $x_2=t^{59}$, $x_3=t^{60}$, and $x_4=t^{67}$
is discussed in Example~4.3 of~\cite{OP},
where they observe that the lattice vectors
corresponding to the minimal binomial generators of $I_{L(1)}$
cannot be arranged as the vertices of the three-dimensional unit cube.
\end{remark}

Throughout this section we use the graded reverse lexicographic order
$\prec$ determined by the grading
$
\deg(x_i)=n_i(m)$, for $1\le i\le4,$
and the variable order
$
x_1\succ x_2\succ x_4\succ x_3.$
Thus, for exponent vectors
$\mathbf a=(a_1,a_2,a_3,a_4)$ and
$\mathbf b=(b_1,b_2,b_3,b_4)$,
we have
\[
{\bf x}^{\mathbf a}\prec {\bf x}^{\mathbf b}
\quad\text{if}\quad
\sum_{i=1}^4 a_i n_i(m)
<
\sum_{i=1}^4 b_i n_i(m),
\]
or if
\[
\sum_{i=1}^4 a_i n_i(m)
=
\sum_{i=1}^4 b_i n_i(m),
\]
and the last nonzero coordinate of
$
(a_1-b_1,\;a_2-b_2,\;a_4-b_4,\;a_3-b_3)$
is positive.


The following notation will be used throughout the remainder of this
section. Let
\[
\begin{aligned}
g_1&=x_1^{m+4}-x_2^{m+1}x_3x_4, \;
g_2=x_4^{m+4}-x_1^{m+1}x_2^3x_3, \;
g_3=x_2^{m+3}-x_1x_3^{m+1}x_4,\\
g_4&=x_1^{m+3}x_2^2-x_3^{m+2}x_4^2, \;
g_5=x_2^{m}x_4^{m+3}-x_1^{m+2}x_3^{m+2}, \;
g_6=x_1^2x_4^{m+2}-x_2x_3^{m+3},\\
g_7&=x_1x_2^{m+2}x_4^{m+1}-x_3^{2m+4}.
\end{aligned}
\]


\begin{theorem}\label{thm:groebner}
For every $m\ge1$, the set
$
G=\{g_1,\ldots,g_7\}$
is the reduced Gr\"obner basis of the lattice ideal $I_{L(m)}$ with
respect to $\prec$. Moreover, $I_{L(m)}$ is generic. The set $G$ is the unique minimal
system of binomial generators of $I_{L(m)}$, and
\[
N(L(m))
=
\{\pm\mathbf u_i\mid1\le i\le7\}.
\]
\end{theorem}

\begin{proof} We first show that $G$ is a Gr\"obner basis for the ideal $J=\langle g_{i} \mid 1 \leq i \leq 7 \rangle$ with
respect to $\prec$. Since every element of $G$ is a binomial
$\mathbf x^{\mathbf u^+}-\mathbf x^{\mathbf u^-}$ with
$\mathbf u=\mathbf u^+-\mathbf u^-\in L(m)$, we have $J\subseteq I_{L(m)}$. The leading monomials of the elements of $G$ are
\[
\operatorname{lm}(g_1)=x_1^{m+4},\;
\operatorname{lm}(g_2)=x_4^{m+4},\;
\operatorname{lm}(g_3)=x_2^{m+3},\;
\operatorname{lm}(g_4)=x_1^{m+3}x_2^2, \]
\[\operatorname{lm}(g_5)=x_2^mx_4^{m+3},\;
\operatorname{lm}(g_6)=x_1^2x_4^{m+2},\;
\operatorname{lm}(g_7)=x_1x_2^{m+2}x_4^{m+1}.
\]

We verify Buchberger's criterion. Since the leading monomials of $g_i$ and $g_j$ are relatively prime for
each of the pairs
\[
\{g_1,g_2\},\;
\{g_1,g_3\},\;
\{g_1,g_5\},\;
\{g_2,g_3\},\;
\{g_2,g_4\},\;
\{g_3,g_6\},
\]
their $S$-polynomials reduce to zero by
Proposition~\ref{prop:coprime}. For the following pairs, we have
\[
\begin{aligned}
S(g_1,g_4)&=x_2^2g_1-x_1g_4=-x_3x_4g_3\stackrel{g_3}{\longrightarrow}0,\\
S(g_1,g_6)&=x_4^{m+2}g_1-x_1^{m+2}g_6=-x_2x_3g_5\stackrel{g_5}{\longrightarrow}0,\\
S(g_2,g_5)&=x_2^mg_2-x_4g_5=-x_1^{m+1}x_3g_3\stackrel{g_3}{\longrightarrow}0,\\
S(g_2,g_6)&=x_1^2g_2-x_4^2g_6=-x_2x_3g_4\stackrel{g_4}{\longrightarrow}0,\\
S(g_3,g_4)&=x_1^{m+3}g_3-x_2^{m+1}g_4=-x_3^{m+1}x_4g_1\stackrel{g_1}{\longrightarrow}0,\\
S(g_3,g_5)&=x_4^{m+3}g_3-x_2^3g_5=-x_1x_3^{m+1}g_2\stackrel{g_2}{\longrightarrow}0,\\
S(g_3,g_7)&=x_1x_4^{m+1}g_3-x_2g_7=-x_3^{m+1}g_6\stackrel{g_6}{\longrightarrow}0,\\
S(g_4,g_6)&=x_4^{m+2}g_4-x_1^{m+1}x_2^2g_6=-x_3^{m+2}g_2\stackrel{g_2}{\longrightarrow}0,\\
S(g_4,g_7)&=x_2^mx_4^{m+1}g_4-x_1^{m+2}g_7=-x_3^{m+2}g_5\stackrel{g_5}{\longrightarrow}0,\\
S(g_5,g_6)&=x_1^2g_5-x_2^mx_4g_6=-x_3^{m+2}g_1\stackrel{g_1}{\longrightarrow}0,\\
S(g_5,g_7)&=x_1x_2^2g_5-x_4^2g_7=-x_3^{m+2}g_4\stackrel{g_4}{\longrightarrow}0,\\
S(g_6,g_7)&=x_2^{m+2}g_6-x_1x_4g_7=-x_3^{m+3}g_3\stackrel{g_3}{\longrightarrow}0.
\end{aligned}
\]

It remains to consider the pairs
$\{g_1,g_7\}$,
$\{g_2,g_7\}$, and
$\{g_4,g_5\}$.

\begin{enumerate}
\item[\textup{(i)}] \emph{The pair $\{g_1,g_7\}$.} Since
$
\operatorname{lcm}\!\left(\operatorname{lm}(g_1),\operatorname{lm}(g_7)\right)
=
x_1^{m+4}x_2^{m+2}x_4^{m+1},$
we have
\[
S(g_1,g_7)
=
x_2^{m+2}x_4^{m+1}g_1
-
x_1^{m+3}g_7.
\]
Hence
\[
\begin{aligned}
S(g_1,g_7)
&=
x_2^{m+2}x_4^{m+1}
(x_1^{m+4}-x_2^{m+1}x_3x_4)
-
x_1^{m+3}
(x_1x_2^{m+2}x_4^{m+1}-x_3^{2m+4})\\
&=
x_1^{m+3}x_3^{2m+4}
-
x_2^{2m+3}x_3x_4^{m+2}.
\end{aligned}
\]
The leading monomial of $S(g_1,g_7)$ is
$x_2^{2m+3}x_3x_4^{m+2}$, which is divisible by
$\operatorname{lm}(g_3)$. We have
\[S(g_1,g_7)
\stackrel{g_3}{\longrightarrow}
x_1^{m+3}x_3^{2m+4}-x_1x_2^mx_3^{m+2}x_4^{m+3}=
-x_1x_3^{m+2}g_5
\stackrel{g_5}{\longrightarrow}
0.\]

\item[\textup{(ii)}] \emph{The pair $\{g_2,g_7\}$.} Since
$
\operatorname{lcm}\!\left(\operatorname{lm}(g_2),\operatorname{lm}(g_7)\right)
=
x_1x_2^{m+2}x_4^{m+4},$
we have
\[
S(g_2,g_7)
=
x_1x_2^{m+2}g_2
-
x_4^3g_7.
\]
Hence
\[
S(g_2,g_7)
=
x_3^{2m+4}x_4^3
-
x_1^{m+2}x_2^{m+5}x_3.
\]
The leading monomial of $S(g_2,g_7)$ is
$x_1^{m+2}x_2^{m+5}x_3$, which is divisible by
$\operatorname{lm}(g_3)$. We have
\[S(g_2,g_7)
\stackrel{g_3}{\longrightarrow}
x_3^{2m+4}x_4^3-x_1^{m+3}x_2^2x_3^{m+2}x_4=
-x_3^{m+2}x_4\,g_4
\stackrel{g_4}{\longrightarrow}
0.\]

\item[\textup{(iii)}] \emph{The pair $\{g_4,g_5\}$.} We distinguish two cases.

\begin{enumerate}
\item[\textup{(a)}] \emph{$m\ge2$.} Since
$
\operatorname{lcm}\!\left(\operatorname{lm}(g_4),\operatorname{lm}(g_5)\right)
=
x_1^{m+3}x_2^mx_4^{m+3},$
we have
\[
S(g_4,g_5)
=
x_2^{m-2}x_4^{m+3}g_4
-
x_1^{m+3}g_5.
\]
Hence
\[
S(g_4,g_5)
=
x_1^{2m+5}x_3^{m+2}
-
x_2^{m-2}x_3^{m+2}x_4^{m+5}.
\]
The leading monomial of $S(g_4,g_5)$ is
$x_1^{2m+5}x_3^{m+2}$, which is divisible by
$\operatorname{lm}(g_1)$. We have
$$S(g_4,g_5)
\stackrel{g_1}{\longrightarrow}
x_1^{m+1}x_2^{m+1}x_3^{m+3}x_4
-
x_2^{m-2}x_3^{m+2}x_4^{m+5}=
-x_2^{m-2}x_3^{m+2}x_4\,g_2
\stackrel{g_2}{\longrightarrow}
0.$$

\item[\textup{(b)}] \emph{$m=1$.} Since
$
\operatorname{lcm}\!\left(\operatorname{lm}(g_4),\operatorname{lm}(g_5)\right)
=
x_1^4x_2^2x_4^4,$
we have
\[
S(g_4,g_5)
=
x_4^4g_4
-
x_1^4x_2g_5.
\]
Hence
\[
S(g_4,g_5)
=
x_1^7x_2x_3^3
-
x_3^3x_4^6.
\]
The leading monomial of $S(g_4,g_5)$ is
$x_1^7x_2x_3^3$, which is divisible by
$\operatorname{lm}(g_1)$. We have
\[S(g_4,g_5)
\stackrel{g_1}{\longrightarrow}
x_1^2x_2^3x_3^4x_4
-
x_3^3x_4^6=
-x_3^3x_4\,g_2
\stackrel{g_2}{\longrightarrow}
0.
\]

\end{enumerate}

\end{enumerate}
Therefore every $S$-polynomial reduces to zero modulo $G$. By
Buchberger's criterion (Theorem~\ref{thm:buchberger}), $G$ is a
Gr\"obner basis of $J$.

Next we show that $G$ is the reduced Gr\"obner basis of $J$ with respect to $\prec$. Since $\operatorname{lc}(g_i)=1$ for every $1\le i\le7$, it remains to
verify that no monomial of an element of $G$ is divisible by the
leading monomial of another element of $G$. This is immediate, since
no leading monomial divides another leading monomial, and every
non-leading monomial contains a positive power of $x_3$, whereas no
leading monomial does. Hence $G$ is the reduced Gr\"obner basis of
$J$. 

Since $x_3$ does not divide any element of $G$, the set $G$ is a Gr\"obner basis of
$J:x_3^{\infty}$ by \cite[Lemma 12.1]{Stu}, thus $J:x_3^{\infty}=J$.

We now show that $J=I_{L(m)}$. It is well known that there is a positive integer $r$ such that $$J=(J:(x_1x_2x_4)^\infty) \cap (J+(x_1x_2x_4)^r).$$ We have \[
x_3^{r(2m+4)}
=
x_2^{(m+1)r}x_4^{mr}\,(x_1x_2x_4)^r
\;-\;
g_7\sum_{i=0}^{r-1}x_3^{(2m+4)(r-1-i)}\bigl(x_1x_2^{m+2}x_4^{m+1}\bigr)^i.
\] Thus $x_3^{r(2m+4)}\in J+(x_1x_2x_4)^r$, and therefore
$\bigl(J+(x_1x_2x_4)^r\bigr):x_3^\infty=S$. Hence
\[
J
=
J:x_3^\infty
=
\Bigl[\bigl(J:(x_1x_2x_4)^\infty\bigr)
\cap
\bigl(J+(x_1x_2x_4)^r\bigr)\Bigr]:x_3^\infty,
\] so \[
J
=
\bigl(J:(x_1x_2x_3x_4)^\infty\bigr)\cap \bigl(J+(x_1x_2x_4)^r\bigr):x_3^\infty
=\bigl(J:(x_1x_2x_3x_4)^\infty\bigr)\cap S=
J:(x_1x_2x_3x_4)^\infty.
\] By \cite[Corollary 2.5]{ES}, we have that $J=I_{L(m)}$. Since each binomial $g_i=\mathbf{x}^{\mathbf{u}_i^+}-\mathbf{x}^{\mathbf{u}_i^-}$
has full support, it follows that $I_{L(m)}$ is generic.

Moreover, no monomial occurring in an element of $G$ divides another
such monomial, and each monomial occurs in exactly one element of $G$.
Hence every monomial occurring in $G$ is indispensable, that is, every
system of binomial generators of $I_{L(m)}$ contains a binomial having
this monomial as one of its terms
\cite[Remark~2.3]{CTV}. Furthermore, the graph associated with $G$ in
the sense of Theorem~3.3 of~\cite{CTV} consists of seven connected
components, each of which is a single edge joining the two monomials
of one of the binomials $g_1,\ldots,g_7$. Therefore every binomial
$g_i$ is indispensable by \cite[Theorem~3.3]{CTV}. Since
$I_{L(m)}$ is generic, \cite[Remark~4.4(3)]{PS}
implies that $G=\{g_1,\ldots,g_7\}$ is the unique minimal system of
binomial generators of $I_{L(m)}$. 

Finally, Proposition~\ref{cor:neighbors} yields
$
N(L(m))=\{\pm\mathbf u_1,\ldots,\pm\mathbf u_7\}.$
\end{proof}

We now determine the $f$-vector of the algebraic Scarf complex of
$I_{L(m)}$. Since $I_{L(m)}$ is generic, \cite[Theorem~4.2]{PS} shows that the algebraic Scarf complex of
$I_{L(m)}$ is the minimal free resolution of $S/I_{L(m)}$.


\begin{proposition}\label{prop:scarf}
For every integer $m\ge1$, the ring $S/I_{L(m)}$ has Betti numbers
$$
\beta_0(S/I_{L(m)})=1, \ \beta_1(S/I_{L(m)})=7, \ \beta_2(S/I_{L(m)})=12, \ \beta_3(S/I_{L(m)})=6,$$ and
$
\beta_i(S/I_{L(m)})=0$ for every $i \geq 4$. Consequently, the algebraic Scarf complex of $I_{L(m)}$ has
$f$-vector
$
(f_0,f_1,f_2)=(7,12,6).$
\end{proposition}
 
\begin{proof} By Proposition~\ref{prop:full-support} and
Theorem~\ref{thm:groebner}, the lattice ideal $I_{L(m)}$ is generic
and has exactly seven minimal binomial generators.

The proof of Corollary~1.3 in~\cite{Ojeda} shows that, for a generic
lattice ideal $I_L$ of codimension~$3$ with $t\ge7$ minimal binomial
generators,
\[
\beta_0(S/I_L)=1, \
\beta_1(S/I_L)=t, \
\beta_2(S/I_L)=2t-2, \
\beta_3(S/I_L)=t-1,
\]
and
$
\beta_i(S/I_L)=0$ for every $i \geq 4$.
Since $I_{L(m)}$ has $t=7$ minimal generators, it follows that
\[
\beta_0(S/I_{L(m)})=1, \
\beta_1(S/I_{L(m)})=7, \
\beta_2(S/I_{L(m)})=12, \
\beta_3(S/I_{L(m)})=6,
\]
and
$
\beta_i(S/I_{L(m)})=0$ for every $i \geq 4$.

By~\cite[Theorem~4.2]{PS}, the minimal free resolution of
$S/I_{L(m)}$ is supported on the algebraic Scarf complex.
Hence the numbers of vertices, edges, and triangles of the algebraic
Scarf complex are
$\beta_1(S/I_{L(m)})$,
$\beta_2(S/I_{L(m)})$, and
$\beta_3(S/I_{L(m)})$, respectively.
Therefore its $f$-vector is $
(f_0,f_1,f_2)=(7,12,6).$
\end{proof}


\section{Restrictions on neighbors of the origin}
\label{sec:lambda}

In this section we study lattice vectors of the form
\[
\mathbf u^{(\mu,\lambda)}
=
\mu\mathbf c_1+\mathbf c_2+\lambda\mathbf c_3,
\]
for positive integers $\mu$ and $\lambda$, where
$\{\mathbf c_1,\mathbf c_2,\mathbf c_3\}$ is an admissible basis of
$L$. Our goal is to determine which such vectors can occur as neighbors
of the origin. We obtain a complete characterization of the vectors
$\mathbf u^{(\mu,\lambda)}$ having positive first and fourth
coordinates that are neighbors of the origin. More precisely, under
conditions~\textup{\textbf{(I)}}, \textup{\textbf{(II)}}, and
\textup{\textbf{(III)}}, we prove that
$\mathbf u^{(\mu,\lambda)}$ is a neighbor of the origin if and only if
$\mu=1$ and $\lambda\le2$.

The case $\mu=1$, namely the vectors
$
\mathbf u^{(1,\lambda)}
=
\mathbf c_1+\mathbf c_2+\lambda\mathbf c_3,
$
was considered in~\cite{OP}, but no complete characterization of the
neighbors of the origin arising from these vectors was obtained. This
motivates our study of the more general family
$\mathbf u^{(\mu,\lambda)}$.

The coordinates of $\mathbf u^{(\mu,\lambda)}$ are
\[
\begin{aligned}
u^{(\mu,\lambda)}_1&=\mu\alpha_{11}-\alpha_{21}-\lambda\alpha_{31},&
u^{(\mu,\lambda)}_2&=-\mu\alpha_{12}-\alpha_{22}+\lambda\alpha_{32},\\
u^{(\mu,\lambda)}_3&=-\mu\alpha_{13}-\alpha_{23}-\lambda\alpha_{33},&
u^{(\mu,\lambda)}_4&=-\mu\alpha_{14}+\alpha_{24}-\lambda\alpha_{34}.
\end{aligned}
\]
Since every $\alpha_{ij}>0$, $u^{(\mu,\lambda)}_3$ is strictly
negative for every positive integers $\mu$ and $\lambda$, so
$(\mathbf u^{(\mu,\lambda)})^+_3=0$.

We say that $\mathbf u^{(\mu,\lambda)}$ has \emph{positive first and
fourth coordinates} if
\[
\mu\alpha_{11}>\alpha_{21}+\lambda\alpha_{31}
\qquad\text{and}\qquad
\alpha_{24}>\mu\alpha_{14}+\lambda\alpha_{34}.
\]

We first derive a general restriction on the parameter $\lambda$ for
arbitrary $\mu$.

\begin{theorem}\label{thm:mu-lambda}
Let $\{\mathbf c_1,\mathbf c_2,\mathbf c_3\}$ be an admissible basis
satisfying condition~\textup{\textbf{(II)}}, and let $\mu\ge1$ and
$\lambda\ge\mu+2$ be integers.
Assume that $\mathbf u^{(\mu,\lambda)}$ has positive first and
fourth coordinates. Then $\mathbf u^{(\mu,\lambda)}$ has full support and is not a
neighbor of the origin in $L$.
\end{theorem}

\begin{proof}
Since $\{\mathbf c_1,\mathbf c_2,\mathbf c_3\}$ is a $\mathbb Z$-basis
of $L$, the vector $\mathbf u^{(\mu,\lambda)}$ belongs to $L$. Condition~\textup{\textbf{(II)}} gives
$\alpha_{32}>\max(\alpha_{12},\alpha_{22})$.
Hence
$2\alpha_{32}>\alpha_{12}+\alpha_{22}$ and
$\alpha_{32}>\alpha_{12}$.
Since $\lambda\ge\mu+2$, we have
\[
\lambda\alpha_{32}
\ge
(\mu+2)\alpha_{32}
=
\mu\alpha_{32}+2\alpha_{32}
>
\mu\alpha_{12}+\alpha_{12}+\alpha_{22}
>
\mu\alpha_{12}+\alpha_{22},
\]
and therefore
$
u^{(\mu,\lambda)}_2
=
\lambda\alpha_{32}
-\mu\alpha_{12}
-\alpha_{22}
>0.
$
By assumption, the first and fourth coordinates of
$\mathbf u^{(\mu,\lambda)}$ are positive. Since the third coordinate is
always negative, it follows that $\mathbf u^{(\mu,\lambda)}$ has full
support. Thus
\[
(\mathbf u^{(\mu,\lambda)})^+
=
\bigl(
\mu\alpha_{11}-\alpha_{21}-\lambda\alpha_{31},
\,
\lambda\alpha_{32}-\mu\alpha_{12}-\alpha_{22},
\,
0,
\,
\alpha_{24}-\mu\alpha_{14}-\lambda\alpha_{34}
\bigr).
\]

It is enough to show that
$\mathbf c_3\in L\setminus\{\mathbf0,\mathbf u^{(\mu,\lambda)}\}$ and
$\mathbf c_3\le(\mathbf u^{(\mu,\lambda)})^+$.

Since
$\mathbf c_3
=
(-\alpha_{31},
\alpha_{32},
-\alpha_{33},
-\alpha_{34})$,
it remains to verify the coordinatewise inequalities.

For the first coordinate,
\[
-\alpha_{31}
<
0
<
\mu\alpha_{11}-\alpha_{21}-\lambda\alpha_{31},
\]
where the second inequality follows from the assumption that the first
coordinate of $\mathbf u^{(\mu,\lambda)}$ is positive.

For the third coordinate, we have
$
-\alpha_{33}
\le
0,$
because the third coordinate of $\mathbf u^{(\mu,\lambda)}$ is
strictly negative.

For the fourth coordinate,
\[
-\alpha_{34}
<
0
<
\alpha_{24}-\mu\alpha_{14}-\lambda\alpha_{34},
\]
where the second inequality follows from the assumption that the fourth
coordinate of $\mathbf u^{(\mu,\lambda)}$ is positive.

For the second coordinate, it is enough to prove
$
(\lambda-1)\alpha_{32}
\ge
\mu\alpha_{12}
+\alpha_{22}.
$
We shall establish the stronger inequality
$
(\lambda-1)\alpha_{32}
>
\mu\alpha_{12}
+\alpha_{22}.
$
Indeed,
\[
(\lambda-1)\alpha_{32}
\ge
(\mu+1)\alpha_{32}
=
\mu\alpha_{32}
+\alpha_{32}
>
\mu\alpha_{12}
+\alpha_{22},
\]
where the last inequality follows from
condition~\textup{\textbf{(II)}}.
Hence
$\mathbf c_3\le(\mathbf u^{(\mu,\lambda)})^+$.

Since
$\{\mathbf c_1,\mathbf c_2,\mathbf c_3\}$ is a basis of $L$,
we have $\mathbf c_3\neq\mathbf0$.
Moreover,
$\mathbf c_3\neq\mathbf u^{(\mu,\lambda)}$
because the first coordinate of $\mathbf c_3$ is negative,
whereas the first coordinate of
$\mathbf u^{(\mu,\lambda)}$ is positive.
Hence
$
\mathbf c_3
\in
L\setminus\{\mathbf0,\mathbf u^{(\mu,\lambda)}\}$ and
$\mathbf c_3
\le
(\mathbf u^{(\mu,\lambda)})^+.$ Lemma~\ref{lem:neighbor} therefore implies that
$\mathbf u^{(\mu,\lambda)}$
is not a neighbor of the origin.
\end{proof}

\begin{remark}
Only the lower bound $\alpha_{32}>\max(\alpha_{12},\alpha_{22})$ in
condition~\textup{\textbf{(II)}} is used in the proof. In particular,
conditions~\textup{\textbf{(I)}} and~\textup{\textbf{(III)}} are not
needed, and the upper bound $\alpha_{32}<\alpha_{12}+\alpha_{22}$
plays no role.
\end{remark}

As an immediate consequence of
Theorem~\ref{thm:mu-lambda}, we obtain the special case $\mu=1$.

\begin{corollary}\label{cor:lambda-general}
Let $\{\mathbf c_1,\mathbf c_2,\mathbf c_3\}$ be an admissible basis
satisfying condition~\textup{\textbf{(II)}}, and let $\lambda\ge3$
be an integer.
If $\mathbf u^{(1,\lambda)}$ has positive first and fourth coordinates,
then $\mathbf u^{(1,\lambda)}$ is not a neighbor of the origin in $L$.
\end{corollary}

The preceding corollary restricts the possible values of $\lambda$ for
neighbors of the form $\mathbf u^{(1,\lambda)}$.

\begin{corollary}\label{cor:lambda}
Let $\{\mathbf c_1,\mathbf c_2,\mathbf c_3\}$ be an admissible basis
with $\alpha_{32}>\max(\alpha_{12},\alpha_{22})$, and let
$\lambda\ge1$ be an integer.
If $\mathbf u^{(1,\lambda)}$ has positive first and fourth coordinates
and is a neighbor of the origin, then $\lambda\le2$.
\end{corollary}

\begin{proof}
Suppose, to the contrary, that $\lambda\ge3$.
Then Corollary~\ref{cor:lambda-general} implies that
$\mathbf u^{(1,\lambda)}$ cannot be a neighbor of the origin,
contradicting the hypothesis.
Therefore $\lambda\le2$.
\end{proof}

Theorem~\ref{thm:mu-lambda} shows that, when $\mu=1$, a necessary
condition for $\mathbf u^{(1,\lambda)}$ to be a neighbor of the origin
is $\lambda\le2$. The next theorem proves that this condition is also
sufficient.

\begin{theorem}\label{thm:converse}
Let $\{\mathbf c_1,\mathbf c_2,\mathbf c_3\}$ be an admissible basis
satisfying conditions~\textup{\textbf{(I)}},
\textup{\textbf{(II)}}, and~\textup{\textbf{(III)}},
and let $\lambda\in\{1,2\}$.
If $\mathbf u^{(1,\lambda)}$ has positive first and fourth
coordinates, then $\mathbf u^{(1,\lambda)}$ is a neighbor of the
origin in $L$.
\end{theorem}

\begin{proof} By Lemma~\ref{lem:neighbor}, it suffices to prove that there is no
$
\mathbf b\in
L\setminus\{\mathbf0,\mathbf u^{(1,\lambda)}\}$
satisfying $
\mathbf b\le(\mathbf u^{(1,\lambda)})^+.$ Assume, to the contrary, that such a vector $\mathbf b$ exists.
Write $\mathbf b=\mu\mathbf c_1+\nu\mathbf c_2+\rho\mathbf c_3$, where $\mu,\nu,\rho\in\mathbb Z$, and define
$p=\mu-1$, $q=\nu-1$, and $r=\rho-\lambda$.
Since $\{\mathbf c_1,\mathbf c_2,\mathbf c_3\}$ is a
$\mathbb Z$-basis of $L$, every element of $L$
has a unique representation with respect to this basis. Hence $\mathbf b=\mathbf u^{(1,\lambda)}$ if and only if $(\mu,\nu,\rho)=(1,1,\lambda)$. By the definitions $p=\mu-1$,
$q=\nu-1$, and 
$r=\rho-\lambda$.
Thus $\mathbf b=\mathbf u^{(1,\lambda)}$ if and only if $(p,q,r)=(0,0,0).$
Similarly, $\mathbf b=\mathbf0$ if and only if $(p,q,r)=(-1,-1,-\lambda)$. 
Since
$
\mathbf b\notin
\{\mathbf0,\mathbf u^{(1,\lambda)}\},$
it follows that
$(p,q,r)\neq(0,0,0)$ and 
$
(p,q,r)\neq(-1,-1,-\lambda)$.

The condition
$
\mathbf b\le(\mathbf u^{(1,\lambda)})^+$
is equivalent to the corresponding coordinatewise inequalities.
Comparing the first, fourth, and second coordinates yields
\begin{align}
p\alpha_{11}&\le q\alpha_{21}+r\alpha_{31},\label{eq:A}\\
q\alpha_{24}&\le p\alpha_{14}+r\alpha_{34},\label{eq:B}\\
r\alpha_{32}&\le p\alpha_{12}+q\alpha_{22}.\label{eq:C}
\end{align}
Since the third coordinate of $\mathbf u^{(1,\lambda)}$ is negative,
the third coordinate of $(\mathbf u^{(1,\lambda)})^+$ is $0$. Therefore the third
coordinate of $\mathbf b$ is nonpositive, that is,
\[
-\mu\alpha_{13}-\nu\alpha_{23}-\rho\alpha_{33}\le0,
\]
which is equivalent to
\begin{equation}\label{eq:D}
(p+1)\alpha_{13}+(q+1)\alpha_{23}+(r+\lambda)\alpha_{33}\ge0.
\end{equation}
We show that the system
\eqref{eq:A}--\eqref{eq:D}
has $(p,q,r)=(0,0,0)$ as its only integer solution by considering all
possible sign patterns of $(p,q,r)$.

\begin{enumerate}
\item[\textbf{Case 1.}] \textit{$p\geq0$, $q \geq 0$, and $r \geq 0$.} We distinguish the following cases.
\begin{enumerate}
\item[(a)] $p=q=0$ and $r>0$. Then \eqref{eq:C} becomes $r\alpha_{32}\le0,$
which is impossible because $r>0$ and $\alpha_{32}>0$.

\item[(b)] $p=0$ and $q>0$. If $r=0$, then \eqref{eq:B} becomes
$
q\alpha_{24}\le0,$
which is impossible since $q>0$ and $\alpha_{24}>0$.

Suppose now that $r>0$.
Then \eqref{eq:B} and \eqref{eq:C} give
$
q\alpha_{24}\le r\alpha_{34}$ and $r\alpha_{32}\le q\alpha_{22}.$
Since $q>0$ and $r>0$, multiplying these inequalities yields
$
\alpha_{24}\alpha_{32}
\le
\alpha_{22}\alpha_{34}.$
On the other hand,
$\alpha_{24}>\alpha_{34}$ by condition~\textup{\textbf{(III)}}
and
$\alpha_{32}>\alpha_{22}$ by condition~\textup{\textbf{(II)}}.
Hence
$
\alpha_{24}\alpha_{32}
>
\alpha_{22}\alpha_{34},$
a contradiction.

\item[(c)] $q=0$ and $p>0$.

If $r=0$, then \eqref{eq:A} becomes
$
p\alpha_{11}\le0,$
which is impossible since $p>0$ and $\alpha_{11}>0$.

Suppose now that $r>0$.
Then \eqref{eq:A} and \eqref{eq:C} give
$
p\alpha_{11}\le r\alpha_{31}$ and $r\alpha_{32}\le p\alpha_{12}$.
Since $p>0$ and $r>0$, multiplying these inequalities yields
$
\alpha_{11}\alpha_{32}
\le
\alpha_{12}\alpha_{31}.$
On the other hand,
$\alpha_{11}>2\alpha_{31}$ by
condition~\textup{\textbf{(I)}}
and
$\alpha_{32}>\alpha_{12}$ by
condition~\textup{\textbf{(II)}}.
Hence
\[
\alpha_{11}\alpha_{32}
>
2\alpha_{31}\alpha_{12}
>
\alpha_{12}\alpha_{31},
\]
a contradiction.

\item[(d)] $p>0$ and $q>0$. From~\eqref{eq:C},
\[
r\le\frac{p\alpha_{12}+q\alpha_{22}}{\alpha_{32}}.
\]
Substituting this upper bound for $r$ into the right-hand sides of
\eqref{eq:A} and \eqref{eq:B} yields
\begin{align}
p(\alpha_{11}\alpha_{32}-\alpha_{12}\alpha_{31})
&\le
q(\alpha_{21}\alpha_{32}+\alpha_{22}\alpha_{31}),
\label{eq:pq1}\\
q(\alpha_{24}\alpha_{32}-\alpha_{22}\alpha_{34})
&\le
p(\alpha_{14}\alpha_{32}+\alpha_{12}\alpha_{34}).
\label{eq:pq2}
\end{align}

By condition~\textup{\textbf{(I)}},
$
\alpha_{11}>\alpha_{21}+2\alpha_{31},$
and by condition~\textup{\textbf{(II)}},
$
\alpha_{32}>\max(\alpha_{12},\alpha_{22}).$
Hence
\begin{align*}
\alpha_{11}\alpha_{32}-\alpha_{12}\alpha_{31}
&>
(\alpha_{21}+2\alpha_{31})\alpha_{32}-\alpha_{12}\alpha_{31} \\
&=
\alpha_{21}\alpha_{32}
+\alpha_{31}(2\alpha_{32}-\alpha_{12}) \\
&>
\alpha_{21}\alpha_{32}
+\alpha_{22}\alpha_{31},
\end{align*}
where the last inequality follows from
$
2\alpha_{32}
>
\alpha_{12}+\alpha_{22}.$

Similarly, condition~\textup{\textbf{(III)}} gives
$
\alpha_{24}>\alpha_{14}+2\alpha_{34},$
and therefore
\begin{align*}
\alpha_{24}\alpha_{32}-\alpha_{22}\alpha_{34}
&>
(\alpha_{14}+2\alpha_{34})\alpha_{32}
-\alpha_{22}\alpha_{34} \\
&=
\alpha_{14}\alpha_{32}
+\alpha_{34}(2\alpha_{32}-\alpha_{22}) \\
&>
\alpha_{14}\alpha_{32}
+\alpha_{12}\alpha_{34},
\end{align*}
where the last inequality again follows from
$
2\alpha_{32}
>
\alpha_{12}+\alpha_{22}.$

Since
\[
\alpha_{11}\alpha_{32}-\alpha_{12}\alpha_{31}
>
\alpha_{21}\alpha_{32}+\alpha_{22}\alpha_{31}>0
\]
and
\[
\alpha_{24}\alpha_{32}-\alpha_{22}\alpha_{34}
>
\alpha_{14}\alpha_{32}+\alpha_{12}\alpha_{34}>0,
\]
we obtain
\[
(\alpha_{11}\alpha_{32}-\alpha_{12}\alpha_{31})
(\alpha_{24}\alpha_{32}-\alpha_{22}\alpha_{34})
>
(\alpha_{21}\alpha_{32}+\alpha_{22}\alpha_{31})
(\alpha_{14}\alpha_{32}+\alpha_{12}\alpha_{34}).
\]
On the other hand, multiplying
\eqref{eq:pq1} and \eqref{eq:pq2}, which is valid since
$p>0$ and $q>0$, gives
\[
(\alpha_{11}\alpha_{32}-\alpha_{12}\alpha_{31})
(\alpha_{24}\alpha_{32}-\alpha_{22}\alpha_{34})
\le
(\alpha_{21}\alpha_{32}+\alpha_{22}\alpha_{31})
(\alpha_{14}\alpha_{32}+\alpha_{12}\alpha_{34}),
\]
a contradiction.

\end{enumerate}

\item[\textbf{Case 2.}] \textit{$p\ge0$, $q\ge0$, and $r<0$.} Since $r<0$ and $\alpha_{31}>0$, inequality~\eqref{eq:A} yields
\[
p\alpha_{11}
\le
q\alpha_{21}+r\alpha_{31}
<
q\alpha_{21}.
\]
As $\alpha_{11}>\alpha_{21}>0$, it follows that $p<q$.

Similarly, since $r<0$ and $\alpha_{34}>0$, inequality~\eqref{eq:B}
gives
\[
q\alpha_{24}
\le
p\alpha_{14}+r\alpha_{34}
<
p\alpha_{14}.
\]
As $\alpha_{24}>\alpha_{14}>0$, we obtain $q<p$, contradicting the inequality $p<q$.

\item[\textbf{Case 3.}] \textit{$p<0$, $q\geq0$, and $r\geq0$.} Since $p<0$ and $\alpha_{12}>0$, inequality~\eqref{eq:C} gives
\[
r\alpha_{32}
\le
p\alpha_{12}+q\alpha_{22}
<
q\alpha_{22}.
\]
As $\alpha_{32}>\alpha_{22}>0$, we obtain $r<q$.

Similarly, since $p<0$ and $\alpha_{14}>0$, inequality~\eqref{eq:B}
gives
\[
q\alpha_{24}
\le
p\alpha_{14}+r\alpha_{34}
<
r\alpha_{34}.
\]
As $\alpha_{24}>\alpha_{34}>0$, we obtain $q<r$, a contradiction.

\item[\textbf{Case 4.}] \textit{$p\geq0$, $q<0$, and $r\geq0$.} Since $q<0$ and $\alpha_{22}>0$, inequality~\eqref{eq:C} gives
\[
r\alpha_{32}
\le
p\alpha_{12}+q\alpha_{22}
<
p\alpha_{12}.
\]
As $\alpha_{32}>\alpha_{12}>0$, we obtain $r<p$.

Similarly, since $q<0$ and $\alpha_{21}>0$, inequality~\eqref{eq:A}
gives
\[
p\alpha_{11}
\le
q\alpha_{21}+r\alpha_{31}
<
r\alpha_{31}.
\]
As $\alpha_{11}>\alpha_{31}>0$, we obtain $p<r$, a contradiction.

\item[\textbf{Case 5.}] \textit{$p<0$, $q<0$, and $r\geq0$.} Since $p<0$, $q<0$, $\alpha_{12}>0$, and $\alpha_{22}>0$, inequality~\eqref{eq:C} gives
\[
r\alpha_{32}
\le
p\alpha_{12}+q\alpha_{22}
<
0.
\]
As $\alpha_{32}>0$, it follows that $r<0$, contradicting the
assumption that $r\ge0$.

\item[\textbf{Case 6.}] \textit{$p<0$, $q\ge0$, and $r<0$.} Since $p<0$, $r<0$, $\alpha_{14}>0$, and $\alpha_{34}>0$, inequality~\eqref{eq:B} gives
\[
q\alpha_{24}
\le
p\alpha_{14}+r\alpha_{34}
<
0.
\]
As $\alpha_{24}>0$, it follows that $q<0$, contradicting the
assumption that $q\ge0$.

\item[\textbf{Case 7.}] \textit{$p\ge0$, $q<0$, and $r<0$.} Since $q<0$, $r<0$, $\alpha_{21}>0$, and $\alpha_{31}>0$, inequality~\eqref{eq:A} gives
\[
p\alpha_{11}
\le
q\alpha_{21}+r\alpha_{31}
<
0.
\]
As $\alpha_{11}>0$, it follows that $p<0$, contradicting the
assumption that $p\ge0$.

\item[\textbf{Case 8.}]  \textit{$p<0$, $q<0$, and $r<0$.} Then $\mu=p+1\leq0$, $\nu=q+1\leq0$, and $\rho=r+\lambda\leq\lambda-1$. We distinguish the following cases.
\begin{enumerate}
\item[(a)] $\lambda=1$. Then $\rho\le0$. Since $\mu\le0$, $\nu\le0$, and $\rho\le0$, while
$\alpha_{13}>0$, $\alpha_{23}>0$, and $\alpha_{33}>0$, each of the
products $\mu\alpha_{13}$, $\nu\alpha_{23}$, and
$\rho\alpha_{33}$ is nonpositive. By~\eqref{eq:D},
\[
\mu\alpha_{13}
+\nu\alpha_{23}
+\rho\alpha_{33}
\ge0.
\]
Since each of the three summands is nonpositive, every summand must
be equal to zero. Therefore,
\[
\mu\alpha_{13}
=
\nu\alpha_{23}
=
\rho\alpha_{33}
=
0.
\]
As $\alpha_{13}$, $\alpha_{23}$, and $\alpha_{33}$ are positive, it
follows that
\[
\mu=\nu=\rho=0.
\]
Hence
$
(p,q,r)=(-1,-1,-1),$
contradicting the fact that
$
(p,q,r)\neq(-1,-1,-\lambda).$

\item[(b)] $\lambda=2$ and $r\le-2$. Then $\rho=r+2\le0.$ Thus
$\mu\le0$, $\nu\le0$, and $\rho\le0$. Arguing exactly as in
case~\textup{(a)}, equation~\eqref{eq:D} again shows that
\[
\mu=\nu=\rho=0.
\]
Hence
$
(p,q,r)=(-1,-1,-2),$
contradicting the assumption that
$
(p,q,r)\neq(-1,-1,-\lambda).$

\item[(c)] $\lambda=2$ and $r=-1$. Then $\rho=r+2=1$. Since $p<0$ and $q<0$, we have
$p\le-1$ and $q\le-1$. Therefore, inequality~\eqref{eq:C} gives
\[
-\alpha_{32}
=
r\alpha_{32}
\le
p\alpha_{12}+q\alpha_{22}
\le
-\alpha_{12}-\alpha_{22}.
\]
Hence
$
\alpha_{12}+\alpha_{22}\le\alpha_{32},$
contradicting the inequality
$
\alpha_{32}<\alpha_{12}+\alpha_{22}$
from condition~\textup{\textbf{(II)}}.
\end{enumerate}
\end{enumerate}

Every possible sign pattern for $(p,q,r)$ leads to a contradiction.
Therefore the system
\eqref{eq:A}--\eqref{eq:D}
has $(p,q,r)=(0,0,0)$ as its only integer solution.
Hence there is no
$
\mathbf b\in
L\setminus
\{\mathbf0,\mathbf u^{(1,\lambda)}\}$
satisfying
$
\mathbf b\le
(\mathbf u^{(1,\lambda)})^+.$
By Lemma~\ref{lem:neighbor},
$\mathbf u^{(1,\lambda)}$
is a neighbor of the origin.
\end{proof}

Combining Corollary~\ref{cor:lambda} with
Theorem~\ref{thm:converse}, we obtain the following characterization.


\begin{corollary}\label{cor:iff}
Let $\{\mathbf c_1,\mathbf c_2,\mathbf c_3\}$ be an admissible basis
satisfying conditions~\textup{\textbf{(I)}},
\textup{\textbf{(II)}}, and~\textup{\textbf{(III)}},
and let $\lambda\ge1$ be an integer.
Assume that $\mathbf u^{(1,\lambda)}$ has positive first and fourth
coordinates. Then $\mathbf u^{(1,\lambda)}$ is a neighbor of the origin
if and only if $\lambda\le2$.
\end{corollary}

\begin{proof}
The necessity follows from Corollary~\ref{cor:lambda}, whereas the
sufficiency is precisely the content of
Theorem~\ref{thm:converse}.
\end{proof}

The preceding characterization concerns the case $\mu=1$.
For arbitrary $\mu$, Theorem~\ref{thm:mu-lambda} yields the following
necessary condition.

\begin{corollary}\label{cor:mu-lambda}
Let $\{\mathbf c_1,\mathbf c_2,\mathbf c_3\}$ be an admissible basis
with $\alpha_{32}>\max(\alpha_{12},\alpha_{22})$, and let $\mu\ge1$
and $\lambda\ge1$ be integers.
If $\mathbf u^{(\mu,\lambda)}$ has positive first and
fourth coordinates and
is a neighbor of the origin, then $\lambda\le\mu+1$.
\end{corollary}

\begin{proof}
If $\lambda\ge\mu+2$ then Theorem~\ref{thm:mu-lambda} gives a
contradiction. Hence $\lambda\le\mu+1$.
\end{proof}

Corollary~\ref{cor:iff} completes the characterization of the
neighbors of the form $\mathbf u^{(1,\lambda)}$ by showing that
$\mathbf u^{(1,\lambda)}$ is a neighbor of the origin if and only if
$\lambda\le2$.
The following theorem shows that the situation changes completely for
$\mu\ge2$: no vector
$\mathbf u^{(\mu,\lambda)}$
having positive first and fourth coordinates is a neighbor of the
origin.

\begin{theorem}\label{thm:mu-never}
Let $\{\mathbf c_1,\mathbf c_2,\mathbf c_3\}$ be an admissible basis
satisfying conditions~\textup{\textbf{(I)}} and~\textup{\textbf{(II)}},
and let $\mu\ge2$ and $\lambda\ge1$ be integers.
If $\mathbf u^{(\mu,\lambda)}$ has positive first and
fourth coordinates,
then $\mathbf u^{(\mu,\lambda)}$ is not a neighbor of the origin.
\end{theorem}

\begin{proof}
By Theorem~\ref{thm:mu-lambda}, we may assume
$\lambda\le\mu+1$. Thus it suffices to consider pairs
$(\mu,\lambda)$ satisfying this inequality.

\begin{enumerate}

\item[\textbf{Case 1.}]
Suppose $(\mu,\lambda)\neq(2,3)$. We shall show that
$\mathbf c_1\in L\setminus\{\mathbf0,\mathbf u^{(\mu,\lambda)}\}$ and
$\mathbf c_1\le(\mathbf u^{(\mu,\lambda)})^+$.
Lemma~\ref{lem:neighbor} will then imply that
$\mathbf u^{(\mu,\lambda)}$ is not a neighbor of the origin.

Since
$\mathbf c_1=(\alpha_{11},-\alpha_{12},-\alpha_{13},-\alpha_{14})$,
it remains to verify the coordinatewise inequalities.

For the first coordinate, we require
$\alpha_{21}\le(\mu-1)\alpha_{11}-\lambda\alpha_{31}$.
By condition~\textup{\textbf{(I)}},
$\alpha_{21}<\alpha_{11}-2\alpha_{31}$.
Hence it is enough to establish
$(\mu-2)\alpha_{11}\ge(\lambda-2)\alpha_{31}$,
and we shall prove the stronger inequality
$(\mu-2)\alpha_{11}>(\lambda-2)\alpha_{31}$.

If $\lambda\le2$, then
$(\lambda-2)\alpha_{31}\le0\le(\mu-2)\alpha_{11}$,
since $\lambda-2\le0$ and $\mu\ge2$.

Assume now that $\lambda\ge3$. If $\mu=2$, then
$\lambda\le\mu+1=3$, so
$(\mu,\lambda)=(2,3)$, contrary to the assumption.
Hence $\mu\ge3$, and therefore
\[
\frac{\lambda-2}{\mu-2}
\le
\frac{\mu-1}{\mu-2}
=
1+\frac1{\mu-2}
\le2.
\]
Condition~\textup{\textbf{(I)}} yields
$\alpha_{11}>\alpha_{21}+2\alpha_{31}>2\alpha_{31}$,
since $\alpha_{21}>0$. Thus
\[\alpha_{11}>
2\alpha_{31}\ge
\frac{\lambda-2}{\mu-2}\alpha_{31},\]
and therefore
$(\mu-2)\alpha_{11}>(\lambda-2)\alpha_{31}$.

For the second coordinate,
$
-\alpha_{12}<0,$
while the second coordinate of $(\mathbf u^{(\mu,\lambda)})^+$ is
nonnegative.

For the third coordinate,
$
-\alpha_{13}<0,$ while the third coordinate of $(\mathbf u^{(\mu,\lambda)})^+$ is equal
to $0$, because the third coordinate of
$\mathbf u^{(\mu,\lambda)}$ is strictly negative.

For the fourth coordinate,
$
-\alpha_{14}<0,$
while the fourth coordinate of $(\mathbf u^{(\mu,\lambda)})^+$ is
positive.

Since
$\{\mathbf c_1,\mathbf c_2,\mathbf c_3\}$ is a basis of $L$,
we have
$\mathbf c_1\neq\mathbf0$.
Moreover,
$\mathbf c_1\neq\mathbf u^{(\mu,\lambda)}$
because the fourth coordinate of $\mathbf c_1$ is negative,
whereas the fourth coordinate of
$\mathbf u^{(\mu,\lambda)}$ is positive.

\item[\textbf{Case 2.}]
Suppose $(\mu,\lambda)=(2,3)$. We shall show that
$\mathbf c_1+\mathbf c_3
\in L\setminus\{\mathbf0,\mathbf u^{(2,3)}\}$ and
$\mathbf c_1+\mathbf c_3
\le(\mathbf u^{(2,3)})^+$.
Lemma~\ref{lem:neighbor} will then imply that
$\mathbf u^{(2,3)}$ is not a neighbor of the origin.

Since
$\mathbf c_1+\mathbf c_3
=
(\alpha_{11}-\alpha_{31},
\alpha_{32}-\alpha_{12},
-\alpha_{13}-\alpha_{33},
-\alpha_{14}-\alpha_{34})$,
we verify the coordinate inequalities.

For the first coordinate,
$\alpha_{11}-\alpha_{31}
\le
2\alpha_{11}-\alpha_{21}-3\alpha_{31}$
is equivalent to
$\alpha_{21}\le\alpha_{11}-2\alpha_{31}$,
which follows immediately from
condition~\textup{\textbf{(I)}}.

For the second coordinate, condition~\textup{\textbf{(II)}} implies that
$\alpha_{12}+\alpha_{22}<2\alpha_{32}$, and hence
$
\alpha_{32}-\alpha_{12}
<
3\alpha_{32}-2\alpha_{12}-\alpha_{22}.$

For the third coordinate,
$
-\alpha_{13}-\alpha_{33}<0,$
while the third coordinate of $(\mathbf u^{(2,3)})^+$ is equal to $0$,
because the third coordinate of $\mathbf u^{(2,3)}$ is strictly
negative.

For the fourth coordinate, $
-\alpha_{14}-\alpha_{34}<0,$
while the fourth coordinate of $(\mathbf u^{(2,3)})^+$ is positive.

Since
$\alpha_{11}>2\alpha_{31}>\alpha_{31}$,
we have
$\mathbf c_1+\mathbf c_3\neq\mathbf0$.
Moreover,
$\mathbf c_1+\mathbf c_3\neq\mathbf u^{(2,3)}$
because the fourth coordinate of
$\mathbf c_1+\mathbf c_3$
is negative,
whereas the fourth coordinate of
$\mathbf u^{(2,3)}$
is positive.
\end{enumerate}

In each case we have constructed a lattice vector
$\mathbf b\in L\setminus\{\mathbf0,\mathbf u^{(\mu,\lambda)}\}$
satisfying
$\mathbf b\le(\mathbf u^{(\mu,\lambda)})^+$.
Lemma~\ref{lem:neighbor} therefore implies that
$\mathbf u^{(\mu,\lambda)}$
is not a neighbor of the origin.
\end{proof}

The preceding results combine to yield a complete characterization of
the vectors $\mathbf u^{(\mu,\lambda)}$ having positive first and
fourth coordinates that are neighbors of the origin.

\begin{corollary}\label{cor:classification}
Let $\{\mathbf c_1,\mathbf c_2,\mathbf c_3\}$ be an admissible basis
satisfying conditions~\textup{\textbf{(I)}},
\textup{\textbf{(II)}}, and~\textup{\textbf{(III)}}.
Let $\mu$ and $\lambda$ be positive integers, and assume that
$\mathbf u^{(\mu,\lambda)}$ has positive first and fourth
coordinates. Then $\mathbf u^{(\mu,\lambda)}$ is a neighbor of the
origin if and only if
$\mu=1$ and 
$\lambda\le2$.
\end{corollary}

\begin{proof}
Suppose first that
$\mathbf u^{(\mu,\lambda)}$
is a neighbor of the origin.
If $\mu\ge2$, then
Theorem~\ref{thm:mu-never}
implies that
$\mathbf u^{(\mu,\lambda)}$
cannot be a neighbor, a contradiction.
Hence $\mu=1$, and
Corollary~\ref{cor:iff}
shows that $\lambda\le2$.

Conversely, if $\mu=1$ and $\lambda\le2$, then
Corollary~\ref{cor:iff}
implies that
$\mathbf u^{(\mu,\lambda)}
=
\mathbf u^{(1,\lambda)}$
is a neighbor of the origin.
\end{proof}

The explicit family of lattices constructed in
Section~\ref{sec:setup} contains neighbors of the origin
with both values $\lambda=1$ and $\lambda=2$,
showing that the bound $\lambda\le2$ can not be improved.

\begin{thebibliography}{99}

\bibitem{BS}
I.~Barany and H.~Scarf, \textit{Matrices with identical sets of neighbors},
Math.\ Oper.\ Res.\ \textbf{23} (1998), no.~4, 863--873.

\bibitem{CTV} H.~Charalambous, A.~Thoma, and M.~Vladoiu, \textit{Binomial fibers and indispensable binomials}, J.\ Symb.\ Comput.\ \textbf{74} (2016), 578--591.

\bibitem{CLO}
D.~Cox, J.~Little, and D.~O'Shea, \textit{Ideals, Varieties, and Algorithms: An Introduction to Computational Algebraic Geometry and Commutative Algebra},
Undergraduate Texts in Mathematics,
Springer, Cham, 2015.

\bibitem{DF}
D.~Dummit and R.~Foote, \textit{Abstract Algebra},
3rd ed.,
John Wiley \& Sons, Hoboken, NJ, 2004.

\bibitem{ES}
D.~Eisenbud and B.~Sturmfels, \textit{Binomial ideals},
Duke Math.\ J.\ \textbf{84} (1996), no.~1, 1--45.

\bibitem{MS}
E.~Miller and B.~Sturmfels, \textit{Combinatorial Commutative Algebra},
Graduate Texts in Mathematics, vol.~227,
Springer, New York, 2005.

\bibitem{Newman}
M.~Newman, \textit{Integral Matrices},
Pure and Applied Mathematics, vol.~45,
Academic Press, New York, 1972.

\bibitem{Ojeda}
I.~Ojeda, \textit{Examples of generic lattice ideals of codimension~3},
Comm.\ Algebra \textbf{36} (2008), no.~1, 279--287.

\bibitem{OP}
I.~Ojeda and P.~Pis\'on-Casares, \textit{On the hull resolution of an affine monomial curve}, J.~Pure Appl.\ Algebra \textbf{192} (2004), 53--67.

\bibitem{PS}
I.~Peeva and B.~Sturmfels, \textit{Generic lattice ideals},
J.~Amer.\ Math.\ Soc.\ \textbf{11} (1998), no.~2, 363--373.

\bibitem{Sabzrou}
H.~Sabzrou, \textit{Scarf lattice ideals},
Portugaliae Math.\ \textbf{68} (2011), no.~4, 369--380.

\bibitem{Stu}
B.~Sturmfels, \textit{Gr\"obner Bases and Convex Polytopes},
University Lecture Series, vol.~8,
Amer.\ Math.\ Soc., Providence, RI, 1996.

\end{thebibliography}

\end{document}
